\documentclass[11pt]{article}
\usepackage{hw-solutions-common}

\begin{document}

\hwsolhead{Written HW 4}{Problem statements from Knight, \emph{Physics for Scientists and
Engineers} (5th ed., Global Edition), Chapter 25. $\dfrac{1}{4\pi\varepsilon_0} =
8.99\times10^{9}$ N$\cdot$m$^2$/C$^2$, $\varepsilon_0 = 8.85\times10^{-12}$
C$^2$/(N$\cdot$m$^2$), and $g = 9.8$ m/s$^2$ throughout.}

%============================================================
\problem{25.6}{2}{What is the electric potential energy of the group of charges in
Figure EX25.6? A $-4.0$ nC charge and a $+3.0$ nC charge sit 3.0 cm apart. A second
$+3.0$ nC charge sits 4.0 cm from the $-4.0$ nC charge, at a right angle to the first
pair, so the two $+3.0$ nC charges are 5.0 cm apart.}

\model{For a set of point charges, the total electric potential energy is the sum of
the potential energies of every distinct pair.}

\visualize{Label the charges $q_1=-4.0$ nC, $q_2=+3.0$ nC, $q_3=+3.0$ nC, with
$r_{12}=3.0$ cm, $r_{13}=4.0$ cm, and $r_{23}=5.0$ cm (the hypotenuse of the 3-4-5 right
triangle formed by the three charges). Three pairs, three terms.}

\solve{
\[
U = \frac{1}{4\pi\varepsilon_0}\left(\frac{q_1q_2}{r_{12}} + \frac{q_1q_3}{r_{13}} +
\frac{q_2q_3}{r_{23}}\right)
\]
\[
= (8.99\times10^{9})\left[\frac{(-4.0\times10^{-9})(3.0\times10^{-9})}{0.030} +
\frac{(-4.0\times10^{-9})(3.0\times10^{-9})}{0.040} +
\frac{(3.0\times10^{-9})(3.0\times10^{-9})}{0.050}\right]
\]
\[
= -3.60\times10^{-6}\text{ J} - 2.70\times10^{-6}\text{ J} + 1.62\times10^{-6}\text{ J}
= \boxed{-4.7\times10^{-6}\text{ J}}.
\]}

\review{The two negative terms (each involving the $-4.0$ nC charge paired with a
$+3.0$ nC charge) outweigh the one positive term (the two $+3.0$ nC charges together),
so the total comes out negative overall --- consistent with the $-4.0$ nC charge being
closer to both $+3.0$ nC charges than they are to each other.}

%============================================================
\problem{25.8}{2}{Three charged particles sit at the corners of an equilateral triangle
with 4.0 cm edges. One particle has charge $+4.0$ nC, a second has charge $+8.0$ nC.
What charge must the third particle have for the total electric potential energy of the
three particles to be zero?}

\model{Every side of an equilateral triangle has the same length, so all three pairwise
separations equal $d=4.0$ cm. Set the sum of the three pairwise potential energies to
zero and solve for the unknown charge $q$.}

\visualize{Call the known charges $q_1=+3.0$ nC and $q_2=+6.0$ nC and the unknown
charge $q$ --- with both given charges positive, their pairwise energy $U_{12}$ is
already positive, so $q$ must be negative to bring the total back down to zero.}

\solve{
\[
U = \frac{1}{4\pi\varepsilon_0 d}\Big[q_1q_2 + q_1q + q_2q\Big] = 0
\;\Rightarrow\;
q_1q_2 = -(q_1+q_2)\,q
\]
\[
(3.0\text{ nC})(6.0\text{ nC}) = -(9.0\text{ nC})\,q
\;\Rightarrow\;
q = \frac{18\text{ nC}^2}{-9.0\text{ nC}} = \boxed{-2.0\text{ nC}}.
\]}

\review{Both given charges are positive, so the pair between them contributes a
positive term no matter what. The only way to zero out the total is a negative third
charge, and that's exactly what falls out --- a good sign the setup is right before
even checking the arithmetic.}

%============================================================
\problem{25.19}{3}{A proton travels along the $x$-axis at points where the electric
potential is $V(x) = (200\text{ V/m})\,x$. As it passes $x=0$ m, the proton's speed is
$2.5\times10^5$ m/s. What speed does it have when it reaches $x=1.0$ m?}

\model{Energy is conserved. The proton's total mechanical energy, kinetic plus
electric potential, is the same at both points.}

\visualize{Call $x=0$ point 1 and $x=1.0$ m point 2, so $V_1=0$ V and $V_2=200$ V. The
proton is positive and moving into higher potential, so it should lose kinetic energy
--- we expect it to slow down.}

\solve{Using $U=qV$ and conservation of energy,
\[
\tfrac12 m v_1^2 + qV_1 = \tfrac12 m v_2^2 + qV_2
\;\Rightarrow\;
v_2 = \sqrt{v_1^2 - \frac{2q}{m}(V_2-V_1)}
\]
\[
= \sqrt{(2.5\times10^5)^2 - \frac{2(1.60\times10^{-19})}{1.67\times10^{-27}}(200\text{
V})} = \boxed{1.5\times10^5\text{ m/s}}.
\]}

\review{The proton does slow down, as expected for a positive charge moving toward
higher potential --- the same rule from the height analogy: a positive charge climbing
toward higher $V$ trades kinetic energy for potential energy, just like a ball climbing
a hill trades kinetic energy for gravitational PE.}

%============================================================
\problem{25.26}{3}{A proton is fired with speed $2.00\times10^5$ m/s from the midpoint
of a parallel-plate capacitor, straight toward the positive plate. The plates are held
at 0 V and $+500$ V. (a) Does the proton reach the positive plate? (b) What speed does
it have when it returns to the negative plate? \emph{Also: do you have enough
information to determine the $x$-position of the proton when it turns around? If so,
what is that point? If not, what information is missing?}}

\model{The potential inside a parallel-plate capacitor rises linearly from the negative
plate to the positive plate, so the midpoint sits exactly halfway between the two plate
voltages: $V_{\text{mid}} = 250$ V. Energy is conserved throughout.}

\visualize{The proton starts at the midpoint moving toward the positive (higher-$V$)
plate, so it decelerates. Whether it reaches the plate depends on whether its initial
kinetic energy is enough to supply the potential-energy increase over the remaining
$500\text{ V}-250\text{ V}=250\text{ V}$ it would still have to climb.}

\solve{(a) The kinetic energy available is
\[
K = \tfrac12 m v^2 = \tfrac12(1.67\times10^{-27}\text{ kg})(2.00\times10^5\text{
m/s})^2 = 3.34\times10^{-17}\text{ J}.
\]
Reaching the positive plate requires climbing the remaining $250$ V from the midpoint,
which costs
\[
\Delta U = e\,\Delta V = (1.60\times10^{-19}\text{ C})(250\text{ V}) =
4.00\times10^{-17}\text{ J}.
\]
Since $K < \Delta U$, \boxed{\text{the proton does not reach the positive plate}} --- it
runs out of kinetic energy first and turns around.

(b) The proton returns through the midpoint with its original speed (energy
conservation is symmetric) and continues on to the negative plate, which is at
\emph{lower} potential than the midpoint, so it speeds up over that stretch:
\[
v_f = \sqrt{v^2 + \frac{2e}{m}(250\text{ V} - 0\text{ V})} = \boxed{2.96\times10^5\text{
m/s}}.
\]

\textbf{Turning point:} Yes, this is determined. At the instant it turns around, all of
its initial kinetic energy has converted to potential energy, so the potential there,
$V_{\text{turn}}$, satisfies $e(V_{\text{turn}}-250\text{ V}) = K = 3.34\times10^{-17}$
J, giving $V_{\text{turn}} = 250\text{ V} + 209\text{ V} = 459$ V. Since $V$ rises
linearly with position, that's $459/500 = 91.8\%$ of the way from the negative plate to
the positive plate --- \emph{as a fraction of the full plate separation}. Converting
that fraction into an actual distance (in mm, say) would need the plate spacing itself,
which isn't needed for anything else in this problem and may or may not be given in the
figure; if it is, multiply the plate spacing by 0.918 to get the position.}

\review{Parts (a) and (b) both hinge on the same idea from lecture: $V$ doesn't have to
change by equal amounts to change the outcome, but here it changes \emph{steadily}, so
distance and potential are directly proportional --- which is exactly what makes the
turning-point fraction (91.8\% of the gap) translate directly into a position, unlike a
capacitor with a bumpier, non-uniform $V(x)$.}

%============================================================
\problem{25.30}{2}{What is the electric potential at the point indicated with the dot
in Figure EX25.30? Three charges $+1.0$ nC, $+2.0$ nC, and $-2.0$ nC sit at the corners
of a symmetric arrangement, each a horizontal distance of 1.5 cm from the dot's vertical
line, at $30^\circ$ from the line to the dot.}

\model{The net potential at a point is the scalar sum of the potentials from each
charge --- no components, no angles needed once each distance to the dot is known.}

\visualize{By the $30^\circ$ geometry, all three charges are the same distance $r$
from the dot: $r = (1.5\text{ cm})/\cos30^\circ = 1.732$ cm.}

\solve{
\[
V = \frac{1}{4\pi\varepsilon_0}\left(\frac{q_1}{r}+\frac{q_2}{r}+\frac{q_3}{r}\right) =
(8.99\times10^9)\left[\frac{1.0\times10^{-9}}{0.01732} +
\frac{2.0\times10^{-9}}{0.01732} + \frac{-2.0\times10^{-9}}{0.01732}\right]
\]
\[
= \boxed{+520\text{ V}}.
\]}

\review{Since $r$ is the same for all three charges, the whole sum collapses to
$\frac{1}{4\pi\varepsilon_0 r}(q_1+q_2+q_3)$ --- the net potential of three equidistant
charges only cares about their \emph{total} charge, $+1.0$ nC in this case. That
shortcut only works because all three distances happen to be equal; it would not apply
if the charges sat at different distances from the dot.}

%============================================================
\problem{25.38}{5}{A $-10.0$ nC point charge and a $+20.0$ nC point charge are 15.0 cm
apart on the $x$-axis, with the $+20.0$ nC charge at the origin. (a) What is the
electric potential at the point on the $x$-axis where the electric field is zero?
(b) At what point on the $x$-axis is the electric potential zero, and what is the
magnitude of the electric field at that point?}

\model{The point where $\vec E=0$ is found from the field magnitudes of the two
charges; the point where $V=0$ is found separately from the (scalar) potentials. These
are two different points in general --- nothing requires them to coincide.}

\visualize{Call $Q_1=+20.0$ nC (at $x=0$) and $Q_2=-10.0$ nC (at $x=15.0$ cm). Since
$Q_1$ is bigger, the field-canceling point must be on the far side of the smaller
charge $Q_2$, at $x>15.0$ cm, where $Q_1$'s field (pointing away from $Q_1$) and
$Q_2$'s field (pointing toward $Q_2$) point in opposite directions and can balance.}

\solve{(a) Setting the two field magnitudes equal at position $x$:
\[
\frac{1}{4\pi\varepsilon_0}\frac{20.0\text{ nC}}{x^2} =
\frac{1}{4\pi\varepsilon_0}\frac{10.0\text{ nC}}{(x-15.0\text{ cm})^2}
\;\Rightarrow\;
x^2 = 2(x-15.0\text{ cm})^2
\]
\[
\Rightarrow x^2-(60.0\text{ cm})x+450\text{ cm}^2=0
\;\Rightarrow\;
x = 51.2\text{ cm (the other root, 8.8 cm, falls between the charges and doesn't fit).}
\]
The potential there is
\[
V = \frac{1}{4\pi\varepsilon_0}\left(\frac{20.0\times10^{-9}}{0.512} +
\frac{-10.0\times10^{-9}}{0.512-0.150}\right) = \boxed{+103\text{ V}}.
\]

(b) Setting the potentials to cancel instead, at a point $x$ between the charges:
\[
\frac{20.0\text{ nC}}{x} + \frac{-10.0\text{ nC}}{15.0\text{ cm}-x} = 0
\;\Rightarrow\;
2(15.0\text{ cm}-x) = x
\;\Rightarrow\;
\boxed{x=10.0\text{ cm}}.
\]
The field there is the vector sum of both contributions (both pointing in $+\hat x$
at this point, since it's to the right of $Q_1$ and to the left of $Q_2$):
\[
E = \frac{1}{4\pi\varepsilon_0}\frac{20.0\times10^{-9}}{(0.100)^2} +
\frac{1}{4\pi\varepsilon_0}\frac{10.0\times10^{-9}}{(0.050)^2} =
\boxed{5.40\times10^4\text{ N/C}}.
\]}

\review{The two answers land on genuinely different points (51.2 cm vs.\ 10.0 cm) ---
a direct illustration that $E=0$ and $V=0$ are unrelated conditions. At the field-zero
point the potential is a nonzero $+103$ V; at the potential-zero point the field is a
very much nonzero $5.4\times10^4$ N/C. Knowing one tells you nothing about the other.}

\end{document}
